% Requested by the Claude supervisor, at Adam Topaz's request, 10 November 2026.
% Mathematical blueprint authored by GPT-6-Astra; no Lean declarations are added.
% Standalone: latexmk -pdf topological_reconstruction_blueprint.tex
\documentclass[11pt]{article}
\usepackage[T1]{fontenc}
\usepackage{lmodern,amsmath,amssymb,amsthm,mathtools,mathrsfs}
\usepackage[margin=1in]{geometry}
\usepackage{microtype,enumitem,booktabs,longtable,array,xcolor,xurl,listings,needspace}
\usepackage[colorlinks=true,linkcolor=blue!45!black,urlcolor=blue!45!black]{hyperref}
\usepackage{fancyhdr}
\pagestyle{fancy}\fancyhf{}
\fancyhead[L]{\small Topological hyperplane reconstruction: M9 blueprint}
\fancyhead[R]{\small\thepage}
\setlength{\headheight}{14pt}
\setlength{\emergencystretch}{3em}
\setlist{itemsep=3pt,topsep=5pt}
\lstset{basicstyle=\ttfamily\small,breaklines=true,columns=fullflexible,keepspaces=true,frame=single}
\newtheorem{theorem}{Theorem}[section]
\newtheorem{lemma}[theorem]{Lemma}
\newtheorem{proposition}[theorem]{Proposition}
\newtheorem{definition}[theorem]{Definition}
\newtheorem{corollary}[theorem]{Corollary}
\theoremstyle{remark}\newtheorem{remark}[theorem]{Remark}
\newcommand{\A}{\mathbb A}
\newcommand{\F}{\mathbb F}
\newcommand{\Z}{\mathbb Z}
\newcommand{\Q}{\mathbb Q}
\newcommand{\OS}{\operatorname{OS}}
\newcommand{\Aut}{\operatorname{Aut}}
\newcommand{\Gal}{\operatorname{Gal}}
\newcommand{\Fix}{\operatorname{Fix}}
\newcommand{\id}{\operatorname{id}}
\newcommand{\GL}{\operatorname{GL}}
\newcommand{\D}{\mathcal D}
\newcommand{\DG}{\mathcal D^{\mathrm{gr}}}
\newcommand{\code}[1]{{\normalfont\nolinkurl{#1}}}
\newcommand{\node}[3]{\par\Needspace{8\baselineskip}\smallskip\noindent\textbf{Node #1.}\quad
 \textsf{Proposed:}\ \code{#2}\par\noindent
 \textsf{Dependencies:} #3.\par\smallskip}
\title{Topological hyperplane reconstruction\\
\large A Lean blueprint upgrading the M8 isomorphisms}
\author{Prepared for Adam Topaz by GPT-6-Astra}
\date{10 November 2026}
\begin{document}
\maketitle
\begin{abstract}
We put the finite-stage convergence topology on the existing marked
quadratic, graded, and cohomological diagram automorphism groups. The
M8 reconstruction maps become topological group isomorphisms with
discrete coefficient units times the field automorphism group with its
pointwise topology. For algebraic extensions this is precisely Mathlib's
existing Krull topology, without normality, separability, perfectness,
or algebraic closure assumptions. We prove the exact weaker condition
for this identification and exhibit a supplied cofinal system over
$\F_p(t)/\F_p$ for which it fails. Scalar quotients, cofinal restriction,
realization independence, and actual ordinary de Rham reconstruction
receive compatible topological upgrades. All new mathematical steps
are proved here; all proposed Lean declarations remain specifications.
\end{abstract}
\tableofcontents
\newpage

\section{Scope, hypotheses, and checked inputs}
\label{sec:scope}

This is the M9 blueprint requested in issue \#202. Source references were
checked against repository commit
\code{56153841b49fa57f217c2dd41e6c39f555f9475f} and Mathlib commit
\code{6db5db0d2279846cc8fb25901bd53221d37a77d8}.
Source inspection is not a new Lean build or kernel replay. The existing
architecture remains authoritative for Lean names and module boundaries.
This document changes neither the M8 mathematical statements nor the
earlier blueprint files. Every name attached to an M9 node is proposed.

Let $k_0\to k$ be an algebra map of fields and let
$G=(k\simeq_{k_0}k)$, written $\Gal(k/k_0)$ even when the extension is
not Galois. Let $C$ be a supplied
\texttt{HyperplaneGT.CofinalSystem k0 k n}. Its stages are finite
$G$-stable arrangements, and every finite set of hyperplanes is contained
in some selected stage. In particular $C$ is directed; its selected
stages need not be closed under unions. Let $R$ be a commutative ring,
with its universe independent of those of both fields, and $S\subset k_0$.
Write $\D_R(C,S)$ and $\DG_R(C,S)$ for the existing
\texttt{DiagramAut R C S} and \texttt{GradedDiagramAut R C S}.

The \emph{classification hypotheses} mean
\[
 \begin{gathered}
  2\leq n,\qquad 0\in S,\qquad \operatorname{Subfield.closure}(S)=k_0,\\
  R\text{ is a nonzero local ring or an integral domain}.
 \end{gathered}
\tag{H}\label{eq:H}
\]
These are exactly the approved M8 hypotheses. Basic topology, adapter
comparisons, and literal restriction maps will not require~\eqref{eq:H}.
The main Krull reconstruction statements additionally assume
\[
                     k/k_0\text{ is algebraic}.                 \tag{A}\label{eq:A}
\]
In Lean this is \texttt{Algebra.IsAlgebraic k0 k}; over fields it supplies
the integrality used by some Mathlib topology lemmas. No normality or
separability is added. Section~\ref{sec:krull} gives the exact weaker
condition and proves that merely supplying $C$ does not replace it.

The coefficient units $R^\times$ are discrete. This is a topology on the
unit group, with no preexisting topology on $R$ required. The product in
the final theorem is the ordinary product of this discrete topology
and Mathlib's Krull topology. In particular, if $R=\Z_\ell$ is used as
an abstract ring, this blueprint gives its units the \emph{discrete}
topology. An adic-coefficient variant is out of scope: it would require
different stage topologies and separate continuity and inverse-continuity
proofs for coefficient extraction and linear automorphisms. Nothing here
asserts that the resulting discrete-coefficient diagram topology is adic.

\subsection{Verified repository references}

Paths below are relative to \code{OrlikSolomon/HyperplaneGT/}; names
are in \code{OrlikSolomon.HyperplaneGT} unless a namespace is stated.
These are existing inputs, not proposed declarations.
\small
\begin{longtable}{@{}p{.42\textwidth}p{.53\textwidth}@{}}
\toprule Source & Existing declaration or fact\\\midrule\endhead
\code{Geometry/Galois/System.lean:71}&
\code{CofinalSystem}; \code{exists_stage_containing} at 99,
\code{exists_ge_arrangement} at 106, \code{exists_ge_finset} at 118.\\
\code{Geometry/Coordinates/Equations.lean:121}&
\code{exists_normal_coefficient_ne_zero_of_isEquation};
\code{zeroLocus_coefficientForm_eq_iff} at 128.\\
\code{Geometry/Galois/CoordinateAction.lean:242}&
\code{affineCoefficients_fieldAffineForm};
\code{isEquation_fieldAffineForm} at 268.\\
\code{Geometry/Galois/Orbits.lean:43}&
\code{finite_scalarOrbit}; \code{fullCofinalSystem} at 172,
the latter under algebraicity.\\
\code{Diagram/Basic/Quadratic.lean:39}&
\code{QuadraticAutFamily}; \code{quadraticDiagramAutSubgroup} at 46;
\code{DiagramAut} at 109; stage/target component API follows.\\
\code{Diagram/Basic/Graded.lean:38}&
\code{GradedAutFamily}; \code{gradedDiagramAutSubgroup} at 44;
\code{GradedDiagramAut} at 110.\\
\code{Algebra/DegreeOne.lean:50}&
\code{adjoin_range_osGenerator}; the canonical degree-one basis
and its inclusion naturality are in the same file.\\
\code{Diagram/Basic/Restriction.lean:198}&
\code{GradedDiagramAut.toQuadratic}; stage/target formulas at
252/256 and injectivity at 262.\\
\code{Reconstruction/LocalRing/Quadratic.lean:104}&
\code{quadraticAutEquiv}; inverse formula at 130;
stage generator formula at 195, target scalar formula at 205.\\
\code{Reconstruction/LocalRing/Graded.lean:41}&
\code{gradedAutEquiv}; \code{gradedClassificationEquiv} at 82;
generator and whole-degree action formulas at 110/121/133.\\
\code{Reconstruction/LocalRing/Quotient.lean:121}&
\code{DiagramAut.scalarQuotientEquiv}; graded version at 130;
inverse section formulas at 154/161.\\
\code{Reconstruction/Domain/Descent.lean:201}&
\code{coefficientDomainAutEquiv}; stage and target formulas
at 313/323.\\
\code{Reconstruction/Domain/Graded.lean:44}&
\code{coefficientDomainGradedAutEquiv}; product classification
\code{coefficientDomainGradedClassificationEquiv} at 88.\\
\code{Reconstruction/Domain/Quotient.lean:84}&
\code{DiagramAut.domainScalarQuotientEquiv}; graded version at 94.\\
\code{Diagram/Cofinal.lean:60}&
\code{DiagramAut.restrictStages}; stage formula at 72 and target
formula at 78; graded restriction at 102.\\
\code{Reconstruction/Cofinal.lean:53}&
\code{cofinalRestrictionEquiv}; graded at 108;
domain and domain-graded versions at 199/254.\\
\code{Cohomology/Adapter/Comparison.lean:198}&
\code{Cohomology.degreeOneDiagramEquiv}; stage/target conjugation
at 205/215 and inverse formulas at 316/326.\\
\code{Cohomology/Adapter/Comparison.lean:290}&
\code{Cohomology.gradedDiagramEquiv}; component formulas at
295/305 and inverse formulas at 334/342.\\
\code{Cohomology/Main/Classification.lean:43}&
\code{Cohomology.DiagramAut.scalarQuotientComparison}; graded
version at 58. Scalar ranges, normality, and adapter compatibility
are in \code{Cohomology/Main/Actions.lean}.\\
\code{Cohomology/Main/Classification.lean:77}&
\code{localQuadraticReconstruction}; \code{localGradedReconstruction}
at 143, \code{localRestrictionEquiv} at 208; domain versions at
374/440/509; actual scalar quotients at 123/189/420/490.\\
\code{Cohomology/Main/Cofinal.lean:162}&
Actual \code{cofinalRestrictionEquiv}; graded at 219,
domain and domain-graded at 303/361.\\
\code{Cohomology/Main/Realization.lean:286}&
\code{realizationQuadraticAutEquiv}; graded at 290;
scalar quotient comparisons at 385/487.\\
\code{Cohomology/DeRham/Diagram.lean:200}&
\code{quadraticTwistEquiv}; \code{gradedTwistEquiv} at 396;
componentwise conjugation and inverse formulas accompany both.\\
\code{Cohomology/DeRham/Reconstruction.lean:78}&
\code{quadraticReconstruction}, \code{gradedReconstruction},
\code{restrictionEquiv}; the final full-system wrappers use
the already constructed characteristic-zero theory.\\
\bottomrule
\end{longtable}
\normalsize

\subsection{Verified Mathlib topology boundary}

The file \code{Mathlib/FieldTheory/KrullTopology.lean} defines
\code{finiteExts} at line 66, \code{galGroupBasis} at 101,
the global \code{krullTopology} instance at 134, and its topological-group
instance immediately after. The theorem
\code{krullTopology_mem_nhds_one_iff} at 144 says exactly
\[
 U\in\mathcal N_G(1)\ \Longleftrightarrow\
 \exists E\subset k\text{ intermediate},\ [E:k_0]<\infty,
                       \ \Fix(E)\subset U.                    \tag{K0}\label{eq:K0}
\]
Here \code{IntermediateField.fixingSubgroup_isOpen} is at 170;
\code{krullTopology_t2} at 186 and the total-separation instance at 224
assume integrality. Crucially, \emph{finite-dimensional} does not mean
finitely generated by possibly transcendental elements. We never install
another global topology on $G$ under the name Krull topology.

\code{Mathlib/Topology/Algebra/ContinuousMonoidHom.lean:310} defines
\code{ContinuousMulEquiv}, notation $\simeq_{\!t}^{\!*}$ in Lean, extending both
\code{MulEquiv} and \code{Homeomorph}.
\code{Mathlib/Topology/Algebra/Group/Quotient.lean} supplies the quotient
topology (29), \code{QuotientGroup.isQuotientMap_mk} (38),
\code{continuous_mk} (42), \code{isOpenQuotientMap_mk} (53), and the
quotient topological-group instance for normal subgroups (151).
Use \code{IsQuotientMap.continuous_iff} in
\code{Mathlib/Topology/Maps/Basic.lean:358} for its universal property.
\code{Mathlib/Topology/Algebra/Group/Basic.lean:283} supplies
\code{isTopologicalGroup_induced}. Products, subtypes, discrete spaces,
and compact products use Mathlib's standard topology instances.
\code{WithTopology} is defined in
\code{Mathlib/Topology/Defs/Basic.lean:250}; its supporting API is in
\code{Mathlib/Topology/WithTopology.lean}.
The existing abbreviation \code{WithDiscreteTopology} in
\code{Mathlib/Topology/Order.lean:588} is also available; wrapping a
factor does not by itself assert an inherited group instance.

\section{The topology of a diagram}
\label{sec:diagram}

\node{T01}{DiagramAut.stageTopology}{M8 A03}
\begin{definition}[Finite-stage convergence]\label{def:t01}
For a stage $A$, put $V_A=\bigoplus_{H\in A}R e_H$ and
$V_T=R e_0\oplus R e_1$. Give each group $\GL_R(V_A)$ and
$\GL_R(V_T)$ the discrete topology, and their product
\[
       P_2=\left(\prod_{A\in C}\GL_R(V_A)\right)\times\GL_R(V_T)
\]
the product topology. Give $\D_R(C,S)$ its existing subgroup's
subspace topology in $P_2$.

For the graded diagram, give every full algebra automorphism group
$\Aut_{R\text{-alg}}(\OS_R(A))$ and the corresponding common-target
group the discrete topology. Give $\DG_R(C,S)$ the subspace topology
in their product. The degree-preservation condition remains part of
the existing subgroup, not an extra coordinate. Define the same topology
on the actual cohomological diagrams using their actual degree-one
modules, respectively their actual diagonal algebras; for ordinary
de Rham use its actual $H^1$, respectively its ordinary graded algebra.
Always include the single original common-target component.
\end{definition}

Thus convergence means eventual equality on any specified finite list
of complete stage automorphisms and on the common target. We use nets
or neighborhoods, not a sequential-convergence definition. There is no
countability assumption on the set of stages.

\paragraph{Lean topology instances.}
The existing family and diagram types are transparent product and subgroup
synonyms. Define the factor topologies explicitly as the discrete topology
$\bot$, then take product and induced topologies; alternatively use
\code{WithTopology} for each discrete factor and pull back along the
component tuple. Expose the chosen diagram topology through a named
definition and a deliberate scoped instance. Use a named/scoped discrete
topology on $R^\times$ as well. Prove component continuity and the
universal property for these chosen topologies. Do not add a global
discrete topology on every \code{LinearEquiv} or \code{AlgEquiv}:
the latter would interfere with the already installed Krull topology.
Do not infer a stage topology from an unrelated topology on $R$.
This instance plumbing must not change a carrier, a group law, or a
field of any existing structure.

\node{T02}{DiagramAut.nhds_one_stageBasis}{T01; M8 G06}
\begin{lemma}[Group structure, separation, and neighborhoods]\label{lem:t02}
All the diagram topologies in T01 make their groups Hausdorff, totally
separated topological groups. For the quadratic diagram, the subgroups
\[
 K_A=\{g:g_T=1,\ g_A=1\}\qquad(A\in C)
                                                               \tag{B}\label{eq:B}
\]
form a neighborhood basis at $1$. The analogous assertion holds for the
graded diagram and for the actual theory and de Rham diagrams.
\end{lemma}
\begin{proof}
Every discrete group is a topological group, as are products and
subgroups. Two different diagram elements differ at a component; its
discrete coordinate separates them by complementary clopen cylinder
sets. This proves Hausdorffness and total separation without
classification or a nonzero-coefficient assumption.

An identity neighborhood in the product contains the simultaneous
identity conditions on finitely many stage coordinates, with the target
condition added if necessary. Choose a common upper stage $B$ by
\code{CofinalSystem.exists_ge_finset}. In the quadratic case, naturality
and injectivity of the finite-support inclusion show that $g_B=1$
implies $g_A=1$ for every selected $A\leq B$. This proves~\eqref{eq:B}.
In the graded case restrict to degree one, apply the same injectivity,
and then use generation of each stage algebra by its degree-one
generators. The actual diagrams have the same properties through their
already proved generator comparisons. Each $K_A$ is a kernel of the
pair of component homomorphisms, hence a normal open subgroup, and its
cosets are clopen. Only the existence of an upper stage was used.
\end{proof}

\node{T03}{DiagramAut.isClosed_embedding}{T01--T02; M8 A01, A03}
\begin{proposition}[Closed subgroup and finite coefficients]\label{prop:t03}
The quadratic and graded diagrams are closed subgroups of their
respective ambient products for every commutative ring $R$.
If $R$ is finite, both are compact and therefore profinite, meaning
compact, Hausdorff, and totally disconnected.
This assertion requires neither classification nor algebraicity.
\end{proposition}
\begin{proof}
For each stage, preservation of its entire quadratic relation span is
a condition on one discrete factor, hence a closed cylinder. For each
inclusion arrow, its naturality equation is a condition on two discrete
factors; any subset of their finite product is clopen, so the equation
defines a closed cylinder. Each marking square similarly depends on
one stage and the common target. The subgroup is the intersection of
all these closed sets, hence closed. There is no finiteness claim about
the number of equations or markings. For the graded diagram, preserving
each degree is again a condition on one discrete coordinate, and each
algebra naturality equation depends on two coordinates. The same
closed-intersection argument applies.

If $R$ is finite, every $V_A$ and $V_T$ is a finite set. Each linear
automorphism group embeds in its finite function set and is finite
discrete, hence compact. The product is compact, and its closed subgroup
is compact. For the graded group, first replace each discrete full
algebra-automorphism factor by its subgroup preserving every degree.
Restriction to degree one injects this subgroup into $\GL_R(V_A)$,
because the algebra is generated by degree one; the same holds at the
target. These discrete subgroups are therefore finite. Their product
is compact, and the diagram's remaining naturality conditions cut out
a closed subspace. T02 gives the separation assertions. In particular
no unproved inverse-limit description or finiteness of an entire
unrestricted algebra-automorphism group is needed.
\end{proof}

\node{T04}{GradedDiagramAut.isEmbedding_toQuadratic}{T01--T02; M8 A01--A03}
\begin{lemma}[The degree-one topology determines the graded topology]\label{lem:t04}
The existing homomorphism
$r:\DG_R(C,S)\to\D_R(C,S)$ is a topological embedding, for arbitrary
commutative coefficients and all dimensions. Whenever the existing
M8 restriction map is bijective, its existing inverse is continuous.
\end{lemma}
\begin{proof}
Restriction at a fixed stage or at the target is a map between discrete
automorphism groups (restricted to the degree-preserving subgroup), so
each coordinate of $r$ is continuous. Conversely, for two graded
diagram elements $g,h$ and any stage $A$, their full stage algebra maps
are equal if and only if their degree-one restrictions agree: the
forward implication is immediate; the reverse follows because the
algebra is generated by the classes $e_H$. The same statement holds
at the common target. Thus the basic equality neighborhoods defining
the graded topology are exactly the pullbacks of the corresponding
quadratic equality neighborhoods. This proves that $r$ is inducing.
Its already proved injectivity makes it an embedding. If $r$ is
bijective, this equality of topologies proves continuity of its inverse.
No assertion that every quadratic automorphism extends is made here;
the existing M8 classification supplies that bijectivity when used.
\end{proof}

\section{Pointwise convergence and the Krull topology}
\label{sec:krull}

We temporarily write $G_{\mathrm{pt}}$ for the same abstract group $G$
with the topology induced by the map $G\to\prod_{a\in k}k$,
$\sigma\mapsto(\sigma(a))_a$, where every copy of $k$ is discrete.
This notation does not replace Mathlib's global Krull instance. Its
identity neighborhoods have the basis $\Fix(B)$ for finite $B\subset k$.

\node{T05}{hyperplaneStabilizer_eq_fixingCoefficients}{M8 G01, G04, G06}
\begin{lemma}[Finite equations and coordinate fibers]\label{lem:t05}
The pointwise topology makes $G$ a Hausdorff totally separated
topological group.
For every hyperplane $H$ there is a finite set $b(H)\subset k$ such that
\[
                \sigma H=H\quad\Longleftrightarrow\quad
                \sigma\text{ fixes every element of }b(H).       \tag{N}\label{eq:N}
\]
If $n\geq1$, put $H_a=\{x:x_0=a\}$. Then
$\sigma H_a=H_{\sigma a}$ and $H_a=H_b$ if and only if $a=b$.
Consequently the topology on $G$ generated by pointwise stabilizers of
finite sets of hyperplanes, or by kernels of its permutations of the
stages of $C$, is precisely $G_{\mathrm{pt}}$.
\end{lemma}
\begin{proof}
To determine $(\sigma\tau)(a)$ near $(\sigma_0,\tau_0)$, require
$\tau(a)=\tau_0(a)$ and $\sigma(\tau_0(a))=\sigma_0(\tau_0(a))$.
To determine $\sigma^{-1}(a)$ near $\sigma_0$, require
$\sigma(\sigma_0^{-1}(a))=a$. These finite evaluation conditions
prove multiplication and inversion continuous. Evaluation separates
different automorphisms by clopen sets.

Choose an equation $c+\sum_i a_i x_i=0$ for $H$, choose an index $j$
with $a_j\ne0$, and divide every coefficient by $a_j$. Let $b(H)$ be
the finite set of these normalized coefficients, including the constant.
The coefficient of $x_j$ is now $1$. The transformed equation has the
same pivot coefficient $1$. Two nonconstant affine equations have the
same zero locus exactly when their coefficient vectors are nonzero
scalar multiples; comparing the pivot forces that scalar to be $1$.
This proves~\eqref{eq:N}, using the existing equation and coordinate
action lemmas. No algebraicity was used.

The equation for $H_a$ gives its transformation rule. Evaluating at
the vector with coordinate $0$ equal to $a$ and other coordinates zero
proves $H_a=H_b\Rightarrow a=b$. Thus every finite field-evaluation
test is a finite hyperplane test. Conversely every finite hyperplane
test is a finite coefficient test by~\eqref{eq:N}. Cofinality places
any finite collection of these hyperplanes into some stage. A stage's
permutation kernel is exactly the pointwise stabilizer of its finite
set of hyperplanes. These two containments prove equality of the
neighborhood bases. The fibers $H_a$ used here are arbitrary geometric
hyperplanes, not only the marking fibers indexed by $S$; the markings
were needed in M8 to identify $G$ over $k_0$.
\end{proof}

\node{T06}{krull_eq_pointwise_iff}{T05; Mathlib Krull topology}
\begin{theorem}[Exact Krull comparison]\label{thm:t06}
For an arbitrary extension $k/k_0$, its Krull topology agrees with
$G_{\mathrm{pt}}$ if and only if
\[
 \text{for every }a\in k\text{ there is }E\subset k,\quad
 [E:k_0]<\infty,\quad \Fix(E)\subset\Fix(\{a\}).             \tag{K}\label{eq:K}
\]
In particular~\eqref{eq:A} implies~\eqref{eq:K}. Under algebraicity,
for every finite collection $A$ of hyperplanes the field
\[
        E_A=k_0\bigl(b(H):H\in A\bigr)
\]
is finite dimensional and
\[
          \Fix(A)=\Gal(k/E_A).                                \tag{F}\label{eq:F}
\]
For every finite intermediate field $E/k_0$ there is a stage $A\in C$
such that $\Fix(A)\subset\Gal(k/E)$, provided $n\geq1$.
\end{theorem}
\begin{proof}
A finite-dimensional intermediate field has a finite vector-space
basis. An automorphism fixes that basis if and only if it fixes the
entire intermediate field, since it fixes $k_0$ and respects addition
and multiplication. Hence every basic Krull subgroup is open in
$G_{\mathrm{pt}}$. Conversely, if~\eqref{eq:K} holds, the stabilizer
of each element contains a basic Krull identity neighborhood and is
therefore open. Finite intersections give continuity of every finite
pointwise test. If the topologies agree, each element stabilizer is a
Krull neighborhood of $1$, and~\eqref{eq:K} follows from~\eqref{eq:K0}.

Under algebraicity take $E=k_0(a)$ to obtain~\eqref{eq:K}. For $E_A$,
adjoining one algebraic element gives a finite extension; repeated
adjoining and the tower law prove finiteness for the finite set of
normalized coefficients. An automorphism fixes the generated field
exactly when it fixes its generators; T05 proves~\eqref{eq:F}.
For the last assertion choose a finite $k_0$-basis $a_1,\ldots,a_m$ of
$E$, and use cofinality to place $H_{a_1},\ldots,H_{a_m}$ in one
stage $A$. Fixing every hyperplane of $A$ fixes every $a_i$ and hence
$E$. Notice that containment of fixing subgroups is the required
statement: it is unnecessary, and generally false, that the chosen
stage have coefficient field exactly $E$.
\end{proof}

Thus algebraicity is a convenient sufficient hypothesis, not a claimed
necessary one: an extension with trivial automorphism group satisfies
\eqref{eq:K} regardless of algebraicity. The most general product
upgrade below may use~\eqref{eq:K} as an explicit proposition; the
recommended public wrappers assume~\eqref{eq:A}. This introduces no
field into \code{CofinalSystem} or a cohomology structure.

\node{T07}{galois_pointwise_compact_of_finite_orbits}{T05}
\begin{proposition}[Compactness from finite orbits]\label{prop:t07}
If every element of $k$ has finite $G$-orbit, $G_{\mathrm{pt}}$ is a
profinite topological group. A supplied $C$ and $n\geq1$ imply this
finite-orbit condition. Algebraicity also implies it directly.
\end{proposition}
\begin{proof}
The topological group and separation statements follow from T05.
Write $O_a=G\cdot a$ and let $X=\prod_{a\in k}O_a$, a product of
finite discrete spaces. Interpret a point of $X$ as a function
$f:k\to k$ satisfying $f(a)\in O_a$. The equations saying that $f$
preserves $0,1,+,\cdot$ and fixes each image of an element of $k_0$
are closed conditions in $X$: each equation depends on finitely many
discrete coordinates. Let $Y$ be their closed intersection. Every
automorphism gives a point of $Y$. Conversely, $f\in Y$ is a unital
field homomorphism, hence injective. If $b\in O_a$ then $O_b=O_a$,
so $f$ restricts to an injective self-map of the finite set $O_a$,
hence a surjection on that set. As $a\in O_a$, every element of $k$
is in the image of $f$. Thus $f$ is an automorphism fixing $k_0$.
The evaluation embedding identifies $G_{\mathrm{pt}}$ homeomorphically
with $Y$. It is therefore compact and totally separated Hausdorff.

For a supplied $C$, choose a finite stable stage containing $H_a$.
Its finiteness and the injectivity of $b\mapsto H_b$ show that $O_a$
is finite. Under algebraicity, the orbit is contained in the finite
root set of the nonzero minimal polynomial of $a$ over $k_0$.
\end{proof}

This proof covers nonnormal and inseparable algebraic extensions.
Alternatively, after the classification upgrade, T03 with coefficients
$\F_2$, $S=k_0$, and the given $n\geq2$ implies compactness of $G$
by the product projection. Neither argument invokes an unproved
compactness theorem for arbitrary algebraic extensions. In particular,
Mathlib's stronger-hypothesis \code{IsGalois} compactness instance need
not be used or made a hypothesis of the main theorem.

\node{T08}{ratFunc_krull_stage_counterexample}{T02, T05--T06; M8 S03; Mathlib rational functions}
\begin{proposition}[A supplied system does not imply the Krull comparison]\label{prop:t08}
Let $k_0=\F_p$, $k=\F_p(t)$, and $n=2$, with $p$ prime. There is a
cofinal system of finite $G$-stable arrangements, $G$ is finite and
nontrivial, and its Krull topology is indiscrete. Its pointwise and
stage permutation topologies are discrete. With $R=\F_2$ and
$S=\{0\}$ all the M8 classification hypotheses hold, but the M8
classification is not a homeomorphism for Mathlib's Krull topology.
\end{proposition}
\begin{proof}
Every nonconstant rational function in $t$ is transcendental over
$\F_p$ (existing \code{RatFunc.transcendental_of_ne_C},
\code{Mathlib/FieldTheory/RatFunc/IntermediateField.lean:120}).
Therefore the only finite-dimensional intermediate field is $\F_p$
itself; all Krull basic identity subgroups equal $G$.

An automorphism is determined by $f=\sigma(t)$. Surjectivity gives
$\F_p(f)=\F_p(t)$. Write $f=P(t)/Q(t)$ in the canonical coprime
numerator/denominator representation. The existing degree theorem
\code{RatFunc.finrank_eq_max_natDegree} in that same file at 165 gives
\[
             1=[\F_p(t):\F_p(f)]=\max(\deg P,\deg Q).
\]
Only finitely many such pairs of polynomials exist over the finite
field, so $G$ is finite. Substitution $t\mapsto t+1$ defines a field
automorphism, with inverse $t\mapsto t-1$, and is nonidentity; this
also works when $p=2$. For completeness the substitution is defined
first on polynomials and then on fractions, where nonzero denominators
stay nonzero because polynomial translation is invertible.

The saturation of any finite set of hyperplanes by the finite group
$G$ is finite and stable. Taking all finite stable arrangements
therefore supplies $C$ directly; do not invoke the existing
\code{fullCofinalSystem} wrapper, whose current hypotheses require
algebraicity. In the pointwise topology the stabilizer of $t$ is
$\{1\}$, so the topology is discrete. Equivalently a stage containing
$H_t$ has trivial permutation kernel. The automorphism group is
nontrivial, hence its indiscrete Krull topology differs from this
discrete topology. The prime field is generated as a subfield by
$\{0\}$, and $\F_2^\times$ is trivial. By M8 the quadratic diagram
has a nontrivial finite Hausdorff, hence discrete, topology. Its
classification with indiscrete $G$ cannot be a homeomorphism; indeed
the forward Galois action is not continuous.
\end{proof}

\section{Upgrading the existing classification maps}
\label{sec:classification}

\node{T09}{quadraticAutContinuousEquiv}{T01--T02, T05--T06; M8 S03, S05}
\begin{theorem}[Quadratic reconstruction as a topological isomorphism]\label{thm:t09}
Assume~\eqref{eq:H}. The existing local-ring map
\code{quadraticAutEquiv}, respectively the domain map
\code{coefficientDomainAutEquiv}, is a topological group isomorphism
\[
       R^\times_{\mathrm{disc}}\times G_{\mathrm{pt}}
                         \ \simeq_{\mathrm{top.grp}}\ \D_R(C,S).
\]
Under~\eqref{eq:K}, and in particular under~\eqref{eq:A}, replace
$G_{\mathrm{pt}}$ by $G$ with its existing Krull topology. This
upgrades the exact existing \code{MulEquiv} in each case.
\end{theorem}
\begin{proof}
Write $\Phi$ for the existing M8 equivalence. Its established formulas
are
\[
 \Phi(u,\sigma)_A(e_H)=u e_{\sigma H},\qquad
 \Phi(u,\sigma)_T(x)=u x.                                    \tag{C}\label{eq:C}
\]
The coefficient ring is nontrivial under~\eqref{eq:H}. The target
component equals the identity if and only if $u=1$, by applying it to
$e_0$ and reading its coefficient at $0$. With $u=1$, equality of the
stage map with the identity is equivalent to $\sigma H=H$ for every
$H\in A$, by independence of the standard basis. Consequently
\[
                    \Phi^{-1}(K_A)=\{1\}\times\Fix(A).        \tag{P}\label{eq:P}
\]
T02 and T05 show that the two sides of~\eqref{eq:P} range through
cofinal identity-neighborhood bases in their respective groups. Since
$\Phi$ is a group isomorphism, translation of these bases proves
continuity in both directions everywhere. T06 supplies exactly the
Krull identification. Both M8 maps have~\eqref{eq:C}, so the proof
applies verbatim to local and domain coefficients, without repeating
fraction-field descent or incidence reconstruction.
\end{proof}

An alternative inverse-continuity proof may use the existing scalar
extraction formula: the scalar is the $(0,0)$ entry of the target
matrix. Having extracted it, a stage containing the coordinate fibers
of a finite basis of $E$ determines the reconstructed automorphism on
$E$. Equation~\eqref{eq:P} is the shorter proof and retains the target
even for systems whose smallest stages are empty.

\node{T10}{gradedAutContinuousEquiv}{T04, T09; M8 S04, S05}
\begin{corollary}[Graded restriction and reconstruction]\label{cor:t10}
Under~\eqref{eq:H}, upgrade the existing local or domain graded
restriction equivalence to
\[
                 \DG_R(C,S)\simeq_{\mathrm{top.grp}}\D_R(C,S).
\]
This statement needs no additional field-extension hypothesis. The
existing \code{gradedClassificationEquiv}, respectively
\code{coefficientDomainGradedClassificationEquiv}, upgrades to
\[
       R^\times_{\mathrm{disc}}\times G_{\mathrm{pt}}
                            \simeq_{\mathrm{top.grp}}\DG_R(C,S),
\]
and to the Krull version under~\eqref{eq:K} or~\eqref{eq:A}.
\end{corollary}
\begin{proof}
T04 and the existing graded restriction bijectivity give the first
upgrade. Compose its inverse with T09 for the second. M8 identifies
this composite with its existing graded classification, so the
underlying map and inverse are unchanged. Its action on degree $q$
is still $u^q$ times the existing Galois action, including degree zero
and every common-target degree. The naturality and marking equations
are exactly the already proved equations.
\end{proof}

\node{T11}{DiagramAut.scalarQuotientContinuousEquiv}{T07, T09--T10; M8 scalar quotients}
\begin{theorem}[Scalar quotients with their quotient topology]\label{thm:t11}
Under~\eqref{eq:H}, let $Z_R$ denote the existing range of the scalar
action in either diagram group. With the actual quotient-group
topology, the existing local and domain scalar quotient equivalences
upgrade to
\[
             \D_R(C,S)/Z_R\simeq_{\mathrm{top.grp}}G_{\mathrm{pt}},
 \qquad      \DG_R(C,S)/Z_R\simeq_{\mathrm{top.grp}}G_{\mathrm{pt}}.
\]
These are Krull isomorphisms under~\eqref{eq:K}. The scalar subgroups
are closed, and the quotient groups are Hausdorff and profinite.
\end{theorem}
\begin{proof}
Under T09 or T10 the scalar range is exactly
$R^\times\times\{1\}$, a closed subgroup since $G_{\mathrm{pt}}$
is Hausdorff. Centrality and normality are already proved in M8.
The extracted Galois coordinate $\pi$ is continuous. The existing
quotient equivalence $\bar\pi$ satisfies $\bar\pi\circ q=\pi$,
where $q$ is the quotient projection. The universal property of the
quotient topology proves $\bar\pi$ continuous. Its already proved
inverse is $\sigma\mapsto q(\Phi(1,\sigma))$, a composite of
continuous maps. The same proof applies in the graded case. T07
gives profiniteness, and T06 supplies the Krull conclusion. This
proves the topology is the quotient topology; it is not defined by
transporting a topology along the quotient equivalence. No equality
between the scalar subgroup and the entire group center is asserted.
\end{proof}

\node{T12}{cofinalRestrictionContinuousEquiv}{T02, T04; M8 cofinal restriction}
\begin{theorem}[Independence of the selected cofinal stages]\label{thm:t12}
Let $C'$ be a supplied cofinal subsystem of $C$. Literal stage
restriction is continuous over any commutative $R$. Under~\eqref{eq:H}
each existing local/domain quadratic/graded cofinal restriction
\code{MulEquiv} upgrades to a topological group isomorphism. No
algebraicity is needed. The induced scalar quotient equivalences
are topological isomorphisms as well.
\end{theorem}
\begin{proof}
Restriction selects the indicated coordinates and leaves the target
coordinate unchanged, so it is continuous by the product universal
property. For inverse continuity in the quadratic case, given an old
stage $A\in C$, choose $B\in C'$ whose arrangement contains $A$.
For two old diagrams $g,h$ agreeing at $B$, inclusion naturality and
injectivity of $V_A\to V_B$ imply $g_A=h_A$. Thus agreement at one
new stage controls the required old stage; the common target is
unchanged. The same argument for finitely many old stages, or T02,
proves the inverse continuous. In the graded case use T04 and the
existing commutative square with degree-one restriction. Bijectivity
is imported from M8, not concluded merely from cofinality. Since the
map and inverse preserve the scalar subgroups, the quotient universal
property proves both quotient directions continuous. For two cofinal
systems without an inclusion, compare both with the system whose
selected stages are the union of their selected families; it is again
cofinal. This uses a union of families, not closure of either family
under arrangement unions.
\end{proof}

\node{T13}{DiagramAut.compactSpace_iff_finite_units}{T03, T07, T09--T11}
\begin{corollary}[Separation and compactness in the classified range]\label{cor:t13}
Under~\eqref{eq:H}, both diagram groups are locally compact and
totally separated Hausdorff. They are compact, equivalently profinite,
if and only if $R^\times$ is finite. In particular finite coefficients
give profinite groups, and $R=\Z$ also gives profinite groups.
Their scalar quotients are always profinite. These conclusions use
the diagram topology even if~\eqref{eq:K} fails; under~\eqref{eq:K}
the quotient is the profinite Krull group.
\end{corollary}
\begin{proof}
By T07 and T09--T10 the groups are products of a discrete group and
a compact totally separated Hausdorff group. Each slice
$\{u\}\times G_{\mathrm{pt}}$ is a compact open neighborhood, proving
local compactness. If the unit group is finite the whole product is
compact. Conversely its continuous surjective projection to the
discrete unit group forces that group to be compact, hence finite:
its singleton open cover must have a finite subcover. For $\Z$ its
units are $\{1,-1\}$. Quotient assertions follow from T11.
For both pure groups with finite coefficients, T03 is the stronger
statement outside the classified range. No compactness claim for
arbitrary infinite-unit coefficient rings is made.
\end{proof}

\section{Actual cohomology and ordinary de Rham}
\label{sec:cohomology}

All theory parameters below are precisely those of the existing
\code{Geometric.OrdinaryTheory} and the M8 adapter: an admissible
smooth-scheme domain, its coefficient algebra and grading, and C1--C5.
No topological-cohomology axiom, continuity field, generation axiom,
or semilinear pullback is added. The topologies concern automorphism
groups of the already constructed finite-stage objects.

\node{T14}{Cohomology.degreeOneDiagramContinuousEquiv}{T01; M8 A04}
\begin{theorem}[The canonical adapters are homeomorphisms]\label{thm:t14}
For any supplied ordinary theory $T$ and arbitrary commutative
coefficients, upgrade its existing comparisons to
\[
 \D_T(C,S)\simeq_{\mathrm{top.grp}}\D_R(C,S),\qquad
 \DG_T(C,S)\simeq_{\mathrm{top.grp}}\DG_R(C,S).
\]
The actual common targets are included. The induced scalar quotient
comparisons are homeomorphisms for the quotient topologies. No
classification, dimension bound, or algebraicity hypothesis is needed.
\end{theorem}
\begin{proof}
Each component of the existing comparison is conjugation by the
canonical degree-one linear equivalence, or the canonical diagonal
algebra equivalence. It is a bijection between discrete factor groups,
hence a homeomorphism. Its inverse is the opposite conjugation.
Taking the product and restricting to the already identified diagram
subgroups proves continuity in both directions. M8 already proves
that these maps preserve scalar ranges. Descending each direction
through the quotient topology proves the last assertion. This also
shows that the topology defined directly from actual components in
T01 agrees with transport through the adapter, rather than defining
it by that transport in the first place.
\end{proof}

\node{T15}{Cohomology.localQuadraticContinuousReconstruction}{T09--T12, T14; M8 E01--E02}
\begin{theorem}[Topological reconstruction for every supplied theory]\label{thm:t15}
Under~\eqref{eq:H} and~\eqref{eq:A}, upgrade the existing maps
\code{localQuadraticReconstruction}, \code{localGradedReconstruction},
\code{domainQuadraticReconstruction}, and
\code{domainGradedReconstruction} to topological group
isomorphisms from $R^\times_{\mathrm{disc}}\times G_{\mathrm{Krull}}$
to the corresponding actual diagram groups, each under its respective
local-ring or integral-domain hypothesis. Upgrade the actual scalar
quotients under the same respective hypotheses to Krull isomorphisms
as well.

The actual local/domain degree-one restriction equivalences and
cofinal restriction equivalences are topological group isomorphisms
already under~\eqref{eq:H}, without~\eqref{eq:A}. All existing
classification, scalar, marking, and restriction squares still commute.
The Krull conclusions also hold with~\eqref{eq:K} replacing~\eqref{eq:A}.
\end{theorem}
\begin{proof}
Compose T09--T10 with the inverse of T14; these composites are the
definitions of the existing actual reconstruction maps. Apply T11
and T14 to the actual scalar quotients. The actual degree-one
restriction is conjugate via T14 to the pure one of T10, so it and
its inverse are continuous. For cofinal restriction use its existing
square with the adapters and T12. All comparisons reuse the existing
\code{MulEquiv}s, and equality of maps is unchanged by adjoining
continuity proofs. Thus the already proved squares and whole-degree
$u^q$ formulas remain literally the same statements after coercion.
\end{proof}

\node{T16}{Cohomology.realizationQuadraticContinuousEquiv}{T14; M8 E02}
\begin{theorem}[Topological realization independence]\label{thm:t16}
For two supplied eligible theories $T_1,T_2$ with the same coefficient
ring and geometric data $k_0,k,n,C,S$, the existing quadratic and graded realization
comparison equivalences are topological group isomorphisms, as are
their existing scalar quotient comparisons. This requires no
classification or algebraicity and allows the independent cohomology
and weight universes, coefficient algebras and gradings, and admissible
domains already admitted in M8. Under the classification
hypotheses the comparisons preserve the unit and Galois parameters
and commute with the upgraded cofinal maps.
\end{theorem}
\begin{proof}
The existing realization comparison is the adapter for $T_1$ followed
by the inverse adapter for $T_2$. Apply T14 twice. On scalar quotients
apply its quotient version twice. The existing M8 component formulas
identify these composites with the canonical actual realization
comparisons, including the common target. Parameter and restriction
compatibility are the existing commuting squares. There is no assumed
morphism of theories and no new cohomology base-change theorem.
\end{proof}

\node{T17}{Cohomology.DeRham.quadraticContinuousReconstruction}{T13--T15; M8 E03}
\begin{corollary}[Actual ordinary de Rham reconstruction]\label{cor:t17}
Let $k$ have characteristic zero. For a supplied $C$, the marking and
dimension hypotheses of~\eqref{eq:H}, and algebraic $k/k_0$, the
existing ordinary de Rham maps upgrade to
\[
 k^\times_{\mathrm{disc}}\times G_{\mathrm{Krull}}
       \simeq_{\mathrm{top.grp}}\D_{\mathrm{dR}}(C,S),\qquad
 k^\times_{\mathrm{disc}}\times G_{\mathrm{Krull}}
       \simeq_{\mathrm{top.grp}}\DG_{\mathrm{dR}}(C,S).
\]
The existing ordinary degree-one restriction is a topological
isomorphism already without algebraicity. The scalar quotients are
Krull isomorphic to $G$, with the quotient topology. Cofinal restriction
and the canonical full-system specializations have the corresponding
topological upgrades. These are unconditional specializations to
the already constructed de Rham theory, not theorems asking for a
supplied realization.
\end{corollary}
\begin{proof}
Equip the ordinary diagram groups with their direct stage topologies
from T01. The existing \code{quadraticTwistEquiv} and
\code{gradedTwistEquiv} conjugate each actual stage and target by the
canonical untwisting equivalence; their inverses do the opposite.
As in T14 these are componentwise homeomorphisms between discrete
factors. Compose their inverses with T15 instantiated at
\texttt{Geometric.DeRham.ordinaryTheory k}. This is exactly the
definition of the existing \code{quadraticReconstruction} and
\code{gradedReconstruction}. Apply the same comparison to restriction,
cofinal maps, and scalar quotients. At the checked revision, ordinary
de Rham scalar subgroups, quotient equivalences, and cofinal restriction
maps are not separately named. Their new algebraic wrappers are
specified as follows. Let $a_C$ be the appropriate quadratic or graded
twist equivalence. Define $Z_{\mathrm{dR}}$ as the range of the already
defined ordinary scalar action. The existing action comparison gives
$a_C(Z_{\mathrm{dR}})=Z_T$, so normality follows from normality of
$Z_T$. The quotient equivalence is the map induced by $a_C$, followed
by the existing theory-level scalar quotient equivalence. Its forward
and inverse continuities follow from the quotient universal property
as in T14. For $C'\subset C$, the ordinary cofinal equivalence is
$a_{C'}^{-1}\circ r_T\circ a_C$, with inverse
$a_C^{-1}\circ r_T^{-1}\circ a_{C'}$, where $r_T$ is the existing
theory-level cofinal equivalence. The component formulas cancel the
stagewise conjugations, so its forward map is literal stage selection
and its target is unchanged. Both directions are continuous by the
already established comparisons. These wrappers add no classification
argument. The full
family is supplied by the existing \code{fullCofinalSystem} under
algebraicity. Finally the same argument works under~\eqref{eq:K}
for a supplied $C$.
\end{proof}

The action remains $k$-linear on logarithmic generators with the
covariant label $H\mapsto\sigma H$, and scales degree $q$ by $u^q$.
It is not identified with semilinear transport of differential forms.
Since a characteristic-zero field has infinitely many nonzero elements,
$k^\times$ with the chosen discrete topology is infinite; these full
de Rham diagram groups are locally compact but not compact. Their
scalar quotients are profinite under the stated hypotheses.

\section{Implementation order, regressions, and acceptance}
\label{sec:implementation}

\subsection{Proposed module boundaries}

The following paths are proposals relative to
\code{OrlikSolomon/HyperplaneGT/}. The architecture owner may split
large leaves; topology is added above the existing M8 modules and
never imported back into their pure algebraic foundations.
\small
\begin{longtable}{@{}p{.34\textwidth}p{.14\textwidth}p{.46\textwidth}@{}}
\toprule Proposed module & Nodes & Dependencies and order\\\midrule\endhead
\code{Topology/Basic.lean}&T01--T02&
\code{Diagram/Basic}; product, induced, discrete, and group topology.
Fix explicit/scoped topology instances and component APIs first.\\
\code{Topology/Compact.lean}&T03&
Basic; compact products and closed subspaces. No reconstruction import.\\
\code{Topology/Restriction.lean}&T04&
Basic; existing degree-one restriction and generation.\\
\code{Topology/Krull.lean}&T05--T06&
Pure coordinate/Galois geometry, \code{Mathlib.FieldTheory.KrullTopology};
normalized coefficient tests and finite-adjoin facts.\\
\code{Topology/GaloisCompact.lean}&T07&
Krull; product compactness and finite orbits. The general field-group
lemma should be stated without arrangements; supplied-$C$ is a corollary.\\
\code{Topology/Classification.lean}&T09--T10&
Basic, Restriction, Krull; existing local/domain classification leaves;
\code{Mathlib.Topology.Algebra.ContinuousMonoidHom}.\\
\code{Topology/Quotient.lean}&T11&
Classification, GaloisCompact; existing scalar quotients and
\code{Mathlib.Topology.Algebra.Group.Quotient}.\\
\code{Topology/Cofinal.lean}&T12&
Basic, Restriction; existing pure cofinal classification. Keep the
literal-map continuity lemma general; import Mathlib's topological
quotient-group module for the quotient assertion.\\
\code{Topology/Consequences.lean}&T13&
Compact, GaloisCompact, Classification, Quotient.\\
\code{Cohomology/Topology/Adapter.lean}&T01--T02, T14&
Actual theory topologies and adapter homeomorphisms; Topology/Basic
and existing A04 comparisons. The T14 quotient portion follows in Main.\\
\code{Cohomology/Topology/Main.lean}&T14--T16&
Adapter, pure topological classification/quotients/cofinal maps and
existing theory Main leaves, including Classification/Actions for
T14's quotient comparisons and Mathlib's topological quotient-group
module. No de Rham imports.\\
\code{Cohomology/DeRham/Topology.lean}&T01--T02, T17&
Theory topological Main plus existing DeRham Diagram/Reconstruction.
Keep this specialization out of the conditional theory facade.\\
\code{OrlikSolomon/Examples/HyperplaneGT/Topology.lean} (full path)&T08;
regressions&
Final topology consumers and rational functions. Keep the rational-function
counterexample imports out of the library core.\\
\bottomrule
\end{longtable}
\normalsize

Give new facade modules their usual sorted imports in
\code{OrlikSolomon.lean}; preserve the existing separation of pure,
conditional-cohomological, and de Rham facades. No coefficient topology
instance may leak from an example into the core. The optional exact
criterion~\eqref{eq:K} belongs in the Krull leaf as an ordinary proposition,
not an additional typeclass prerequisite of all M8 declarations.
The pure Basic leaf implements only the pure-diagram portions of
T01--T02. Their actual-theory portions are implemented in the adapter
leaf, and their ordinary de Rham portions in its specialization leaf.
These repeated node references specify specializations of the same
statements, not imports from the pure core into a realization layer
in the wrong direction.

Each upgraded map should be built by adjoining
\code{continuous_toFun} and \code{continuous_invFun} proofs to its
existing \code{MulEquiv}. Export a theorem identifying
\code{toMulEquiv} with that exact existing equivalence, plus the
coercion/application formula. Retain the existing inverse and component
formulas rather than restating the classification. Do not mark broad
expansion equations as simp rules where that would interfere with
ordinary two-sided equivalence cancellation. Existing normality
instances for scalar ranges should be reused.

\subsection{Required regressions}

\begin{enumerate}[leftmargin=*]
\item \textbf{Characteristic-two coefficients.} Instantiate the Krull
upgrade at $R=\Z/2$ and an algebraic geometric extension with nontrivial
automorphism group. The scalar factor is trivial, so the actual
quadratic and graded groups are topologically $G$. Check the original
common target and the quotient projection formulas, not just group
cardinalities. Nothing in the proof divides by $2$.
\item \textbf{Nilpotent coefficients.} Use
$R=\F_2[\epsilon]/(\epsilon^2)$, the existing dual-number fixture.
The unit $1+\epsilon\ne1$ is detected on the target in degree one,
although its square is $1$ and its degree-two action is trivial.
Thus the topology and inverse proof cannot replace the target or
degree-one component by degree two alone. The ring is finite local,
so its diagram groups are profinite and the nonidentity scalar is a
closed discrete factor.
\item \textbf{Domain coefficients.} Use $R=\Z$ and check the domain
upgrade and its underlying M8 map. The units $\{1,-1\}$ are finite,
so the diagram is profinite even though the ring is infinite.
At a coefficient field the local/domain topological maps agree
because their underlying M8 maps agree and continuity proofs are
proof-irrelevant.
\item \textbf{Dimension and non-union closure.} Run the final wrappers
at $n=2$. Reuse M8's actual cofinal subsystem with two singleton stages
whose union is excluded and which admits an enlargement through a
third hyperplane. Check that T02 and T12 call the cofinal enlargement
lemmas; no union-closure field or silently replaced system is allowed.
\item \textbf{Finite Galois group.} For any finite algebraic extension,
take the finite intermediate field $E=k$ in~\eqref{eq:K0}. The trivial
subgroup is open, so the Krull topology is discrete; $G$ is finite.
The product and quotient statements specialize accordingly. A
nontrivial quadratic extension tests both the section and quotient.
\item \textbf{Infinite Galois group.} Use an algebraic closure of
$\F_p$ with the full stable system, or a supplied algebraic extension
with infinite $G$. T07 gives compactness; an infinite compact discrete
space cannot exist, so $G$ is not discrete. For the algebraic-closure
example, Frobenius is an automorphism since $x\mapsto x^p$ is
injective and every element has a $p$th root. Its positive powers are
distinct: if its $m$th power were identity then every element would
be a root of $X^{p^m}-X$, forcing the algebraic closure to be finite.
An algebraically closed field is infinite, since for a finite field
the polynomial $1+\prod_{a\in k}(X-a)$ has no root. Thus this is a
genuine infinite example. Pin a stage kernel as an open subgroup;
do not accidentally give the whole field automorphism group the
discrete topology through an \code{AlgEquiv} instance.
\item \textbf{Failure without the Krull condition.} Formalize T08,
including an actual finite-stable cofinal-system witness and the
nontrivial translation automorphism. This checks that the new
algebraicity wrapper hypothesis has a mathematical purpose. It must
not be weakened to ``finite orbits'' while retaining Mathlib's Krull
topology. Use the already proved rational-function degree theorem.
\item \textbf{Unclassified and degenerate cases.} T01--T04 and T14
retain arbitrary commutative coefficients, including the zero ring,
and permit $n=0$. A dimension-zero system may have only empty stages;
the common target must still be present. Do not instantiate the
classified product theorem there. The topological argument evaluates
formal coefficient vectors, not points of arrangement complements,
so it does not impose a rational-point hypothesis. Existing M8
point-free examples remain valid consumers of all applicable maps.
\item \textbf{Actual theories and de Rham.} Reuse the supplied-theory
fixtures for $\F_2$, the dual numbers, and $\Z$; their realization
claims remain conditional on those supplied theories. Separately
instantiate T17 with the constructed characteristic-zero de Rham
theory, retaining ordinary $H^1$, full ordinary cohomology, and its
untwisting square. Test both inverse continuities and all quotient
topologies, not only continuity of the forward actions.
\end{enumerate}

\subsection{Book tracking and completion criteria}

There are seventeen logical M9 nodes T01--T17, each with exactly one
labelled statement in this document. Their dependencies are recorded
next to the statements; M8 S03--S05 and E01--E03 remain formally
proved algebraic inputs. The book should add separate M9 boxes and
mark each as planned until the corresponding topology statement is
formalized. A product theorem with pointwise topology alone does not
complete its Krull wrapper. A continuous bijection without a proof of
inverse continuity does not complete a topological-isomorphism node.

Acceptance requires direct component-topology definitions, topological
group instances, the explicit basis comparison~\eqref{eq:P}, actual
quotient-topology proofs, cofinal and two-realization compatibility,
and the ordinary de Rham specialization. Audit the inferred topology
on $k\simeq_{k_0}k$ against \texttt{krullTopology k0 k}. Audit the
underlying equivalence of every upgraded theorem against the exact
M8 declaration. All new hypotheses must be visible, especially
algebraicity or~\eqref{eq:K}; no assumption may assert a reconstruction
conclusion. No M8 source statement or cohomology structure needs
alteration. The resulting library proves the minimal requested
discrete-coefficient theorem and makes no adic-coefficient claim.

\end{document}
